\documentclass[11pt,a4paper]{article}
\usepackage[utf8]{inputenc}
\usepackage[margin=2.5cm]{geometry}
\usepackage{amsmath,amssymb,amsfonts}
\usepackage{array}
\usepackage{hyperref}

\hypersetup{
    colorlinks=true,
    linkcolor=blue,
    urlcolor=blue,
    citecolor=blue
}

\setlength{\parindent}{0pt}
\setlength{\parskip}{6pt}

\newcommand{\kotakdefinisi}[2]{%
\noindent\fbox{%
\parbox{\dimexpr\textwidth-2\fboxsep-2\fboxrule\relax}{%
\textbf{#1}\\[4pt]
#2
}%
}\par\vspace{6pt}
}

\begin{document}

% --- Header Catatan Kuliah ---
\begin{center}
    {\LARGE \textbf{CATATAN KULIAH MA4171 TEORI KONTROL LINEAR}}\\[6pt]
    {\Large \textbf{Topik 04: Keterkontrolan dan Keterobservasian Sistem}}\\[4pt]
    \textit{Cakupan: Minggu ke-6 sampai Minggu ke-7}\\[2pt]
    \textbf{Referensi Utama:} Katsuhiko Ogata, \textit{Modern Control Engineering} (Bab 9.2, 9.6, \& 9.7)
\end{center}
\vspace{-4pt}
\hrule
\vspace{10pt}

\tableofcontents
\vspace{10pt}
\hrule
\vspace{12pt}

\section{Pendahuluan \& Konsep Struktural Sistem}

Konsep \textbf{Keterkontrolan (\textit{Controllability})} dan \textbf{Keterobservasian (\textit{Observability})} diperkenalkan pertama kali oleh Rudolf E. Kalman pada awal tahun 1960-an. Keduanya merupakan sifat struktural fundamental dalam teori kontrol modern berbasis ruang keadaan (\textit{state-space}):

\begin{itemize}
    \item \textbf{Keterkontrolan} menjawab pertanyaan: \textit{"Dapatkah masukan kontrol $u(t)$ menggerakkan seluruh variabel keadaan internal $x(t)$ dari sembarang titik awal ke titik tujuan yang diinginkan dalam waktu berhingga?"}
    \item \textbf{Keterobservasian} menjawab pertanyaan: \textit{"Dapatkah kondisi awal seluruh variabel keadaan internal $x(0)$ direkonstruksi/diestimasi secara unik hanya dari riwayat pengukuran keluaran $y(t)$ dan masukan $u(t)$?"}
\end{itemize}

Keterkontrolan merupakan syarat mutlak agar kita dapat merancang pengendali umpan balik keadaan (\textit{state feedback / pole placement}) untuk menempatkan kutub-kutub lup tertutup pada posisi sembarang di bidang kompleks. Keterobservasian merupakan syarat mutlak dalam merancang pengamat keadaan (\textit{state observer / estimator}).

\vspace{6pt}

\section{Keterkontrolan Keadaan (\textit{State Controllability})}

Diberikan sistem LTI kontinu:
\begin{equation*}
    \dot{x}(t) = A x(t) + B u(t), \quad x(t) \in \mathbb{R}^n, \; u(t) \in \mathbb{R}^m
\end{equation*}

\kotakdefinisi{Definisi Keterkontrolan Keadaan Lengkap:}{
Sistem $(A, B)$ dikatakan \textbf{terkontrol keadaan lengkap} (\textit{completely state controllable}) jika untuk sembarang keadaan awal $x(t_0) = x_0$ dan sembarang keadaan akhir $x_1$, terdapat sinyal masukan kendali tak-terkendala $u(t)$ yang dapat mentransfer keadaan sistem dari $x_0$ ke $x_1$ dalam interval waktu berhingga $t_0 \le t \le t_1$.
}

\subsection{Kriteria Matriks Keterkontrolan Kalman}
Dari solusi domain waktu sistem:
\begin{equation*}
    x(t_1) = e^{A(t_1 - t_0)} x(t_0) + \int_{t_0}^{t_1} e^{A(t_1 - \tau)} B u(\tau) \, d\tau
\end{equation*}
Berdasarkan \textbf{Teorema Cayley-Hamilton}, matriks eksponensial $e^{At}$ dapat diekspansikan sebagai kombinasi linear hingga orde $n-1$:
\begin{equation*}
    e^{At} = \sum_{k=0}^{n-1} \alpha_k(t) A^k
\end{equation*}
Sehingga integral kendali berada pada rentang kolom dari matriks:
\begin{equation*}
    M_c = \begin{bmatrix} B & AB & A^2 B & \dots & A^{n-1} B \end{bmatrix} \in \mathbb{R}^{n \times nm}
\end{equation*}

\kotakdefinisi{Teorema Keterkontrolan Kalman:}{
Sistem kontinu LTI $(A, B)$ dengan $n$ variabel keadaan bersifat \textbf{terkontrol keadaan lengkap} jika dan hanya jika matriks keterkontrolan $M_c$ memiliki rank penuh (\textit{full row rank}):
\begin{equation*}
    \operatorname{rank}(M_c) = \operatorname{rank} \begin{bmatrix} B & AB & A^2 B & \dots & A^{n-1} B \end{bmatrix} = n
\end{equation*}
Untuk sistem SISO ($m = 1$), $M_c$ berbentuk matriks bujursangkar $n \times n$, sehingga syarat ekuivalennya adalah $\det(M_c) \ne 0$.
}

\subsection{Gramian Keterkontrolan (\textit{Controllability Gramian})}
Keterkontrolan pada selang waktu $[0, t_1]$ dapat juga diuji melalui matriks Gramian keterkontrolan:
\begin{equation*}
    W_c(0, t_1) = \int_0^{t_1} e^{A(t_1 - \tau)} B B^T e^{A^T(t_1 - \tau)} \, d\tau \in \mathbb{R}^{n \times n}
\end{equation*}
Sistem $(A, B)$ terkontrol jika dan hanya jika $W_c(0, t_1)$ bersifat \textbf{simetris dan definit positif} ($W_c > 0$) untuk setiap $t_1 > 0$.

Sinyal kendali dengan energi minimum $\int_0^{t_1} u^T(t) u(t) \, dt$ yang mentransfer $x(0) = x_0$ menuju $x(t_1) = x_1$ diberikan oleh:
\begin{equation*}
    u(t) = B^T e^{A^T(t_1 - t)} W_c^{-1}(0, t_1) \left[ x_1 - e^{At_1} x_0 \right]
\end{equation*}

\subsection{Uji Popov-Belevitch-Hautus (PBH Test) untuk Keterkontrolan}
Uji PBH merupakan uji berbasis nilai eigen yang sangat elegan dan praktis:

\kotakdefinisi{Uji PBH Keterkontrolan:}{
Sistem $(A, B)$ terkontrol keadaan lengkap jika dan hanya jika:
\begin{equation*}
    \operatorname{rank} \begin{bmatrix} sI - A & B \end{bmatrix} = n, \quad \forall s \in \mathbb{C}
\end{equation*}
Dalam praktiknya, kita hanya perlu menguji untuk nilai-nilai $s = \lambda_i$, di mana $\lambda_i$ adalah nilai eigen dari matriks $A$.
}

\subsection{Uji Keterkontrolan Gilbert (Bentuk Diagonal / Jordan)}
Jika matriks $A$ dapat didiagonalkan dengan transformasi keserupaan $x = P z$ sehingga $\Lambda = P^{-1} A P = \operatorname{diag}(\lambda_1, \lambda_2, \dots, \lambda_n)$ dengan nilai eigen berbeda:
\begin{equation*}
    \dot{z}(t) = \Lambda z(t) + \tilde{B} u(t), \quad \tilde{B} = P^{-1} B
\end{equation*}
\textbf{Syarat Keterkontrolan Gilbert}: Sistem $(A, B)$ terkontrol jika dan hanya jika \textbf{tidak ada satu pun baris pada matriks $\tilde{B}$ yang seluruh elemennya bernilai nol}.

\vspace{6pt}

\section{Keterobservasian Keadaan (\textit{State Observability})}

Diberikan sistem LTI kontinu:
\begin{equation*}
    \dot{x}(t) = A x(t) + B u(t), \quad y(t) = C x(t) + D u(t)
\end{equation*}
dengan $x(t) \in \mathbb{R}^n$, $u(t) \in \mathbb{R}^m$, dan $y(t) \in \mathbb{R}^p$.

\kotakdefinisi{Definisi Keterobservasian Keadaan Lengkap:}{
Sistem $(A, C)$ dikatakan \textbf{terobservasi keadaan lengkap} (\textit{completely state observable}) jika untuk sembarang keadaan awal $x(t_0) = x_0$, nilai $x_0$ dapat ditentukan secara unik dari pengukuran keluaran $y(t)$ dan masukan kendali $u(t)$ pada interval waktu berhingga $t_0 \le t \le t_1$.
}

\subsection{Kriteria Matriks Keterobservasian Kalman}
Tanpa mengurangi keumuman (karena efek masukan $u(t)$ diketahui dan dapat dikurangkan), perhatikan respon alami tanpa masukan $u(t) = 0$:
\begin{equation*}
    y(t) = C e^{At} x_0 = \sum_{k=0}^{n-1} \alpha_k(t) C A^k x_0
\end{equation*}
Agar $x_0$ dapat ditentukan secara unik dari fungsi $y(t)$, matriks keterobservasian harus memiliki rank kolom penuh.

\kotakdefinisi{Teorema Keterobservasian Kalman:}{
Sistem kontinu LTI $(A, C)$ dengan $n$ variabel keadaan bersifat \textbf{terobservasi keadaan lengkap} jika dan hanya jika matriks keterobservasian $M_o$ memiliki rank penuh (\textit{full column rank}):
\begin{equation*}
    \operatorname{rank}(M_o) = \operatorname{rank} \begin{bmatrix} C \\ CA \\ CA^2 \\ \vdots \\ CA^{n-1} \end{bmatrix} = n
\end{equation*}
Untuk sistem dengan keluaran skalar ($p = 1$), $M_o$ adalah matriks $n \times n$, sehingga syarat ekuivalennya adalah $\det(M_o) \ne 0$.
}

\subsection{Gramian Keterobservasian (\textit{Observability Gramian})}
Matriks Gramian keterobservasian didefinisikan sebagai:
\begin{equation*}
    W_o(0, t_1) = \int_0^{t_1} e^{A^T \tau} C^T C e^{A \tau} \, d\tau \in \mathbb{R}^{n \times n}
\end{equation*}
Sistem $(A, C)$ terobservasi jika dan hanya jika $W_o(0, t_1)$ bersifat \textbf{simetris dan definit positif} ($W_o > 0$) untuk setiap $t_1 > 0$.

\subsection{Uji PBH untuk Keterobservasian}
\kotakdefinisi{Uji PBH Keterobservasian:}{
Sistem $(A, C)$ terobservasi keadaan lengkap jika dan hanya jika:
\begin{equation*}
    \operatorname{rank} \begin{bmatrix} sI - A \\ C \end{bmatrix} = n, \quad \forall s \in \mathbb{C}
\end{equation*}
Cukup diuji pada seluruh nilai eigen $s = \lambda_i$ dari matriks $A$.
}

\subsection{Uji Keterobservasian Gilbert}
Jika sistem berada dalam bentuk kanonik diagonal $\dot{z} = \Lambda z, y = \tilde{C} z$ dengan $\tilde{C} = C P$ dan nilai-nilai eigen berbeda:\\
\textbf{Syarat Keterobservasian Gilbert}: Sistem $(A, C)$ terobservasi jika dan hanya jika \textbf{tidak ada satu pun kolom pada matriks $\tilde{C}$ yang seluruh elemennya bernilai nol}.

\vspace{6pt}

\section{Prinsip Dualitas Kalman}

Rudolf E. Kalman membuktikan relasi matematis yang sangat indah antara masalah keterkontrolan dan keterobservasian:

\kotakdefinisi{Teorema Dualitas Kalman:}{
Pasangan sistem $(A, B)$ bersifat \textbf{terkontrol lengkap} jika dan hanya jika pasangan dualnya $(A^T, B^T)$ bersifat \textbf{terobservasi lengkap}.
}

\subsection{Tabel Ekuivalensi Sistem Primer dan Sistem Dual}
\begin{center}
\small
\begin{tabular}{|l|c|c|}
\hline
\textbf{Properti} & \textbf{Sistem Primer $\mathcal{S}$} & \textbf{Sistem Dual $\mathcal{S}^*$} \\ \hline
\textbf{Persamaan Ruang Keadaan} & $\dot{x} = Ax + Bu, \; y = Cx$ & $\dot{z} = A^T z + C^T v, \; w = B^T z$ \\ \hline
\textbf{Matriks Sistem} & $A$ & $A^T$ \\ \hline
\textbf{Matriks Masukan / Keluaran} & Masukan: $B$, Keluaran: $C$ & Masukan: $C^T$, Keluaran: $B^T$ \\ \hline
\textbf{Matriks Keterkontrolan} & $M_c = \begin{bmatrix} B & AB & \dots & A^{n-1}B \end{bmatrix}$ & $M_c^* = M_o^T$ \\ \hline
\textbf{Matriks Keterobservasian} & $M_o = \begin{bmatrix} C \\ CA \\ \vdots \\ CA^{n-1} \end{bmatrix}$ & $M_o^* = M_c^T$ \\ \hline
\textbf{Kondisi Rank} & $\operatorname{rank}(M_c) = n$ & $\operatorname{rank}(M_o^*) = \operatorname{rank}(M_c^T) = n$ \\ \hline
\end{tabular}
\end{center}

\textbf{Manfaat Utama Dualitas}: Setiap algoritma atau teorema untuk perancangan kontrol umpan balik keadaan (\textit{state feedback}) dapat dipetakan langsung untuk merancang pengamat keadaan (\textit{state observer}) dengan memanfaatkan transpos matriks.

\vspace{6pt}

\section{Keterkontrolan Keluaran (\textit{Output Controllability})}

Pada beberapa aplikasi praktis, kita hanya tertarik mengontrol variabel keluaran sistem $y(t)$, bukan seluruh variabel keadaan internal $x(t)$.

\kotakdefinisi{Definisi Keterkontrolan Keluaran:}{
Sistem dikatakan \textbf{terkontrol keluaran lengkap} (\textit{completely output controllable}) jika untuk sembarang keluaran awal $y(t_0)$ dan sembarang keluaran akhir $y_1$, terdapat masukan $u(t)$ yang dapat mentransfer keluaran dari $y(t_0)$ ke $y_1$ dalam waktu berhingga $t_1 - t_0$.
}

\kotakdefinisi{Syarat Rank Keterkontrolan Keluaran:}{
Sistem $\dot{x} = Ax + Bu, y = Cx + Du$ dengan $p$ keluaran terkontrol keluaran lengkap jika dan hanya jika:
\begin{equation*}
    \operatorname{rank}(M_{co}) = \operatorname{rank} \begin{bmatrix} CB & CAB & CA^2 B & \dots & CA^{n-1} B & D \end{bmatrix} = p
\end{equation*}
}

\textbf{Hubungan Antara Keterkontrolan Keadaan dan Keterkontrolan Keluaran}:
\begin{itemize}
    \item Jika sistem terkontrol keadaan lengkap dan matriks $C$ memiliki baris-baris yang bebas linear ($\operatorname{rank}(C) = p$), maka sistem \textbf{pasti} terkontrol keluaran lengkap.
    \item Keterkontrolan keluaran \textbf{tidak menjamin} keterkontrolan keadaan sistem (keadaan internal bisa saja tidak terkontrol).
\end{itemize}

\vspace{6pt}

\section{Keterstabilan (\textit{Stabilizability}) \& Keterdeteksian (\textit{Detectability})}

Jika suatu sistem tidak terkontrol atau tidak terobservasi lengkap, kita masih dapat mengendalikan atau mengestimasinya asalkan dinamika tak terkontrol / tak terobservasi tersebut stabil dengan sendirinya:

\begin{itemize}
    \item \textbf{Keterstabilan (\textit{Stabilizability})}: Sistem $(A, B)$ dikatakan \textit{stabilizable} jika seluruh mode tak terkontrolnya bersifat stabil asimtotik ($\operatorname{Re}(\lambda) < 0$).
    \item \textbf{Keterdeteksian (\textit{Detectability})}: Sistem $(A, C)$ dikatakan \textit{detectable} jika seluruh mode tak terobservasinya bersifat stabil asimtotik ($\operatorname{Re}(\lambda) < 0$).
\end{itemize}

\kotakdefinisi{Uji PBH untuk Keterstabilan dan Keterdeteksian:}{
\begin{align*}
    (A, B) \text{ stabilizable} &\iff \operatorname{rank} \begin{bmatrix} sI - A & B \end{bmatrix} = n, \quad \forall s \in \mathbb{C} \text{ dengan } \operatorname{Re}(s) \ge 0 \\
    (A, C) \text{ detectable} &\iff \operatorname{rank} \begin{bmatrix} sI - A \\ C \end{bmatrix} = n, \quad \forall s \in \mathbb{C} \text{ dengan } \operatorname{Re}(s) \ge 0
\end{align*}
}

\vspace{6pt}

\section{Dekomposisi Kanonik Kalman \& Pembatalan Kutub-Nol}

\subsection{Empat Sub-ruang Dekomposisi Kalman}
Melalui transformasi keserupaan yang tepat, sembarang ruang keadaan dapat didekomposisikan menjadi 4 sub-ruang ortogonal:
\begin{equation*}
    x = \begin{bmatrix} x_{co} \\ x_{c\bar{o}} \\ x_{\bar{c}o} \\ x_{\bar{c}\bar{o}} \end{bmatrix}
    \begin{array}{l}
        \leftarrow \text{Terkontrol dan Terobservasi} \\
        \leftarrow \text{Terkontrol dan Tak Terobservasi} \\
        \leftarrow \text{Tak Terkontrol dan Terobservasi} \\
        \leftarrow \text{Tak Terkontrol dan Tak Terobservasi}
    \end{array}
\end{equation*}

\subsection{Kaitan dengan Fungsi Transfer \& Pembatalan Kutub-Nol (\textit{Pole-Zero Cancellation})}
Fungsi transfer sistem $G(s) = C(sI - A)^{-1} B + D$ **hanya merepresentasikan sub-sistem yang terkontrol dan terobservasi ($x_{co}$)**:
\begin{equation*}
    G(s) = C_{co} (sI - A_{co})^{-1} B_{co} + D
\end{equation*}

\kotakdefinisi{Teorema Realisasi Minimal \& Pembatalan Kutub-Nol:}{
\begin{enumerate}
    \item Representasi ruang keadaan $(A, B, C, D)$ berorde $n$ merupakan \textbf{realisasi minimal} dari fungsi transfer $G(s)$ jika dan hanya jika sistem bersifat \textbf{terkontrol lengkap dan terobservasi lengkap}.
    \item Jika terjadi \textbf{pembatalan kutub-pembuat nol (\textit{pole-zero cancellation})} pada fungsi transfer $G(s)$, maka representasi ruang keadaan yang dibangun pasti \textbf{kehilangan sifat keterkontrolan, keterobservasian, atau keduanya}.
\end{enumerate}
}

\vspace{6pt}

\section{Keterkontrolan \& Keterobservasian Sistem Diskrit}

Untuk sistem LTI waktu diskrit:
\begin{equation*}
    x(k+1) = G x(k) + H u(k), \quad y(k) = C x(k) + D u(k), \quad x(k) \in \mathbb{R}^n
\end{equation*}

\kotakdefinisi{Matriks Keterkontrolan dan Keterobservasian Diskrit:}{
\begin{align*}
    M_c &= \begin{bmatrix} H & GH & G^2 H & \dots & G^{n-1} H \end{bmatrix} \in \mathbb{R}^{n \times nm} \\
    M_o &= \begin{bmatrix} C \\ CG \\ CG^2 \\ \vdots \\ CG^{n-1} \end{bmatrix} \in \mathbb{R}^{np \times n}
\end{align*}
Sistem diskrit terkontrol lengkap $\iff \operatorname{rank}(M_c) = n$.\\
Sistem diskrit terobservasi lengkap $\iff \operatorname{rank}(M_o) = n$.
}

\vspace{6pt}

\section{Contoh Soal dan Pembahasan Komprehensif}

\noindent\textbf{Contoh 1: Uji Keterkontrolan dan Keterobservasian Kalman}\\
Diberikan sistem kontinu:
\begin{equation*}
    A = \begin{bmatrix} 1 & 1 \\ 0 & -1 \end{bmatrix}, \quad B = \begin{bmatrix} 1 \\ 0 \end{bmatrix}, \quad C = \begin{bmatrix} 1 & 0 \end{bmatrix}
\end{equation*}
Periksalah keterkontrolan dan keterobservasian sistem tersebut.

\textit{Penyelesaian:}

1. \textbf{Uji Keterkontrolan}:
Hitung $AB$:
\begin{equation*}
    AB = \begin{bmatrix} 1 & 1 \\ 0 & -1 \end{bmatrix} \begin{bmatrix} 1 \\ 0 \end{bmatrix} = \begin{bmatrix} 1 \\ 0 \end{bmatrix}
\end{equation*}
Bentuk matriks keterkontrolan $M_c$:
\begin{equation*}
    M_c = \begin{bmatrix} B & AB \end{bmatrix} = \begin{bmatrix} 1 & 1 \\ 0 & 0 \end{bmatrix}
\end{equation*}
Determinan $\det(M_c) = (1)(0) - (1)(0) = 0 \implies \operatorname{rank}(M_c) = 1 < 2$.  
\textbf{Kesimpulan}: Sistem \textbf{TIDAK terkontrol keadaan lengkap}.

2. \textbf{Uji Keterobservasian}:
Hitung $CA$:
\begin{equation*}
    CA = \begin{bmatrix} 1 & 0 \end{bmatrix} \begin{bmatrix} 1 & 1 \\ 0 & -1 \end{bmatrix} = \begin{bmatrix} 1 & 1 \end{bmatrix}
\end{equation*}
Bentuk matriks keterobservasian $M_o$:
\begin{equation*}
    M_o = \begin{bmatrix} C \\ CA \end{bmatrix} = \begin{bmatrix} 1 & 0 \\ 1 & 1 \end{bmatrix}
\end{equation*}
Determinan $\det(M_o) = (1)(1) - (0)(1) = 1 \ne 0 \implies \operatorname{rank}(M_o) = 2$.  
\textbf{Kesimpulan}: Sistem \textbf{terobservasi keadaan lengkap}.

\vspace{10pt}

\noindent\textbf{Contoh 2: Analisis Pembatalan Pole-Zero}\\
Diberikan fungsi transfer sistem:
\begin{equation*}
    G(s) = \frac{s + 1}{(s + 1)(s + 2)} = \frac{s + 1}{s^2 + 3s + 2}
\end{equation*}
Susun representasi bentuk kanonik terkontrol (\textit{controllable canonical form}), kemudian hitung matriks $M_c$ dan $M_o$.

\textit{Penyelesaian:}

Dari $G(s) = \dfrac{s + 1}{s^2 + 3s + 2}$, koefisien penyebut $a_1 = 3, a_2 = 2$ dan pembilang $b_1 = 1, b_2 = 1$.  
Bentuk kanonik terkontrolnya adalah:
\begin{equation*}
    A = \begin{bmatrix} 0 & 1 \\ -2 & -3 \end{bmatrix}, \quad B = \begin{bmatrix} 0 \\ 1 \end{bmatrix}, \quad C = \begin{bmatrix} 1 & 1 \end{bmatrix}
\end{equation*}

1. \textbf{Matriks Keterkontrolan}:
\begin{equation*}
    AB = \begin{bmatrix} 0 & 1 \\ -2 & -3 \end{bmatrix} \begin{bmatrix} 0 \\ 1 \end{bmatrix} = \begin{bmatrix} 1 \\ -3 \end{bmatrix} \implies M_c = \begin{bmatrix} 0 & 1 \\ 1 & -3 \end{bmatrix}
\end{equation*}
$\det(M_c) = -1 \ne 0 \implies \operatorname{rank}(M_c) = 2$ (\textbf{Terkontrol Lengkap}).

2. \textbf{Matriks Keterobservasian}:
\begin{equation*}
    CA = \begin{bmatrix} 1 & 1 \end{bmatrix} \begin{bmatrix} 0 & 1 \\ -2 & -3 \end{bmatrix} = \begin{bmatrix} -2 & -2 \end{bmatrix} \implies M_o = \begin{bmatrix} 1 & 1 \\ -2 & -2 \end{bmatrix}
\end{equation*}
$\det(M_o) = (1)(-2) - (1)(-2) = 0 \implies \operatorname{rank}(M_o) = 1 < 2$ (\textbf{TIDAK Terobservasi}).

\textbf{Interpretasi Fisis}: Pembatalan faktor $(s + 1)$ menyebabkan mode $\lambda = -1$ menjadi tidak terobservasi pada keluaran.

\vspace{10pt}

\noindent\textbf{Contoh 3: Uji PBH dengan Parameter Tak Diketahui}\\
Diberikan matriks sistem:
\begin{equation*}
    A = \begin{bmatrix} -1 & 0 \\ 0 & -2 \end{bmatrix}, \quad B = \begin{bmatrix} 1 \\ \alpha \end{bmatrix}, \quad C = \begin{bmatrix} \beta & 1 \end{bmatrix}
\end{equation*}
Tentukan nilai parameter $\alpha$ dan $\beta$ agar sistem bersifat terkontrol lengkap dan terobservasi lengkap.

\textit{Penyelesaian:}

Nilai eigen dari matriks diagonal $A$ adalah $\lambda_1 = -1$ dan $\lambda_2 = -2$.

1. \textbf{Uji Keterkontrolan via PBH}:
\begin{itemize}
    \item Untuk $s = -1$:
    \begin{equation*}
        \begin{bmatrix} -1 \cdot I - A & B \end{bmatrix} = \begin{bmatrix} 0 & 0 & 1 \\ 0 & 1 & \alpha \end{bmatrix} \implies \text{rank} = 2 \quad (\forall \alpha)
    \end{equation*}
    \item Untuk $s = -2$:
    \begin{equation*}
        \begin{bmatrix} -2 \cdot I - A & B \end{bmatrix} = \begin{bmatrix} -1 & 0 & 1 \\ 0 & 0 & \alpha \end{bmatrix} \implies \text{rank} = 2 \iff \alpha \ne 0
    \end{equation*}
\end{itemize}
Jadi sistem terkontrol lengkap jika dan hanya jika $\mathbf{\alpha \ne 0}$.

2. \textbf{Uji Keterobservasian via PBH}:
\begin{itemize}
    \item Untuk $s = -1$:
    \begin{equation*}
        \begin{bmatrix} -1 \cdot I - A \\ C \end{bmatrix} = \begin{bmatrix} 0 & 0 \\ 0 & 1 \\ \beta & 1 \end{bmatrix} \implies \text{rank} = 2 \iff \beta \ne 0
    \end{equation*}
    \item Untuk $s = -2$:
    \begin{equation*}
        \begin{bmatrix} -2 \cdot I - A \\ C \end{bmatrix} = \begin{bmatrix} -1 & 0 \\ 0 & 0 \\ \beta & 1 \end{bmatrix} \implies \text{rank} = 2 \quad (\forall \beta)
    \end{equation*}
\end{itemize}
Jadi sistem terobservasi lengkap jika dan hanya jika $\mathbf{\beta \ne 0}$.

\end{document}
